\section{Related Work}
\label{sec:relatedwork}

Simulation is a key technology in DTs, envisioned as a tool that can be adopted across all phases of the PT lifecycle~\cite{Boschert2016}.
%
Despite the strong overlap between simulation and DTs making it difficult to draw a line between the technologies~\cite{WOOLEY2023940}, it is becoming more and more clear that, especially in the IoT domain, DTs extend the capability of ``static'' simulations by incorporating real-time data~\cite{PHANDEN2021174}.

Most DT simulations integrate simulators and IoT data processing with custom architectures and ad-hoc methods, lacking principled integration. The DT Simulation Bridge integrates and orchestrates multiple simulations with different models in a coherent framework.
%
At a surface level, it is then similar to distributed simulation~\cite{2000daglib} and co-simulation~\cite{Gomes&18}. 
%
There are, though, fundamental differences in the approach.
%
Distributed simulation usually concerns the execution of a single complex simulation on distributed nodes.
%
Differently, co-simulation explores the possibility of composing different simulation models of individual parts to simulate a complex system and coordinate the progress of time between the different independent simulation units.
%
The purpose of the DT Simulation Bridge is, instead, to connect one DT with multiple simulations.
%, independent from one another.
%
The coordination and combination of such simulation results are expected to be managed by the DT and not by the bridge itself, which only acts as an intermediary.
%
Nevertheless, despite these differences in purposes, there are some significant aspects from these two related areas which influence the design of the bridge. 

Integrating simulations requires standardising the interaction across components.
%
For distributed simulation, the High-Level Architecture (HLA) defines a set of constraints to design interoperable federations of simulations which can be executed on a distributed runtime infrastructure~\cite{568652}.
%
Despite in our case simulations do not need to exchange data between each other, the need for a standardized interface is essential to allow the DT to seamlessly interact with multiple different simulations.
%
Similarly, in the co-simulation domain, the Functional Mockup Interface (FMI) is an industrial de facto standard \cite{FMIStandard2.0} to bridge different simulation models.
%
There, models of different nature may declare connection points to create input/output chains which combine the output of individual solvers.
%
We follow this principle in the design of the DT-SB to support simulation models which can be updated over time with the data that is captured by the DT itself. 